\documentclass[10pt,a4paper]{amsart}
\usepackage[margin=19mm]{geometry}
\usepackage{amsmath,amssymb,mathtools}
\usepackage[scr=rsfs,cal=euler]{mathalfa}
\usepackage{xfrac}
\usepackage[unicode,hidelinks,bookmarks=false]{hyperref}
\newcommand{\ie}{i.e.\ }
\newcommand{\cf}{cf.\ }
\newcommand{\resp}{respectively\ }
\newcommand{\Subsection}{\subsection}
\providecommand{\nobreakdash}{\nobreak\hbox{-}}
\setlength{\emergencystretch}{2em}
\multlinegap0pt
\usepackage{amssymb}
\usepackage{mathtools}
\usepackage{xcolor}
\usepackage{tikz}
\usepackage{dsfont}
\usepackage[shortlabels]{enumitem}
\newtheorem{claim}{Claim}
\title{Tropical Tevelev Degrees}
\author{Renzo Cavalieri, Erin Dawson}
\date{Source: arXiv:2407.20025v1 (CC BY 4.0)}
\begin{document}
\maketitle

\noindent\emph{Source and licence.} Tropical Tevelev Degrees. Version arXiv:2407.20025v1 (CC BY 4.0), distributed by arXiv under the Creative Commons Attribution 4.0 International licence, http://creativecommons.org/licenses/by/4.0/, as recorded in the arXiv licence metadata for this exact version and shown on the abstract page https://arxiv.org/abs/2407.20025v1. Full licence notice: \texttt{licenses/arxiv-2407.20025-CC-BY-4.0.txt}.\par\smallskip

\noindent\textit{Editorial scope.} Counted unit: the article's claim claim (source0.tex lines 2673--2675). The article's complete original proof of it is delivered with it (source0.tex lines 2677--2697, one \texttt{proof} environment of 21 lines). No mathematics is written, repaired or re-derived: the statement and the proof are the article's own lines, in the article's own notation, and no retained line is substituted or re-flowed.  No further source-local context is needed: every cross-reference of every retained line resolves inside the two delivered spans.  Nothing was omitted from any retained span, so the package carries no omission record.  No retained line contains a citation, so the wrapper carries no bibliography and no bibliography entry.  No retained line contains a figure environment, an image-inclusion command or drawing code, so no image is delivered and the package declares no asset dependency.  Editorial formatting only, in addition: an AMS \texttt{amsart} preamble that restates the article's own declarations and macros used by the retained lines, namely the environments \texttt{claim} on separate counters, together with the standard packages those lines need.  The retained environments print in this document as Claim 1 in order of appearance, whereas the article prints the same items with the numbering of its own full document.  The complete native source of the article as distributed in the arXiv e-print for this exact version is bundled unchanged as \texttt{sources/163041/source0.tex}.
\begin{claim}
    The determinant of the block corresponding to the genus part of the cover is equal to $2^g$. 
\end{claim}

\begin{proof}
    
The block that corresponds to the genus part of the cover is itself made of smaller blocks.    We assign  the following matrices, $M_U$ and $M_D$, respectively, to the genus fragments $U$ and $D$. %both of which have determinant 2.
\[ M_U=\left[ \begin{array}{cc}
1 & 0 \\
d & 2
\end{array} \right]
%
M_D=\left[ \begin{array}{cc}
1 & 0 \\
d-1 & 2
\end{array} \right]
\]
The genus block of the matrix
is constructed as a sequence of $1\times 1$ and $2\times 2$ diagonal blocks as follows. 
There is an initial diagonal entry of $2$ corresponding to the first loop.
%to right along our graphs. %The first loop contributes a single 2 to the matrix. 
Then we have a block diagonal %matrix containing blocks that look like
entry of $M_U$ or $M_D$ depending on which type of genus fragment has been used to form the second genus. Between every loop there is an edge that is part of the active path and thus there is a 1 on the diagonal between any two consecutive $M_U$/$M_D$ blocks.
In conclusion, we have $g-1$ $(2\times 2)$-blocks of determinant $2$, one diagonal entry of $2$ and $g-1$ more diagonal entries of $1$, and the claim follows.
\end{proof}

\end{document}
